\chapter{Verification and Validation}

Verification and Validation is a whole life-cycle process that must be applied at each stage in the software process.

It has two principal objectives:
\vspace{-0.5\baselineskip}
\begin{itemize}[parsep=0.25em]
   \item The \orange{discovery of defects} in a system;
   \item The \orange{assessment} of whether or not the system is \orange{useful} and \orange{useable} in an operational situation.
\end{itemize}


\section{Verification}
\ding{70} \textit{Are we building the product right?}

Verification is the process of checking that a software achieves its
goal without any bugs. It verifies whether the developed software fulfills requirements, without necessarily executing the code.

\textbf{Activities involved}: Code reviews, inspections, walkthroughs, static analysis, unit testing.

\section{Validation}
\ding{70} \textit{Are we building the right product?}

Validation is the process of checking that the software meets the needs of the customer and fulfills its high level requirements. It ensures that the right product is built, and it is done by executing the code.

\textbf{Activities involved}: System testing, acceptance testing, performance testing.

\section{V and V Techniques}